\documentclass[12pt,a4paper]{article}
\usepackage[UTF8]{ctex}
\usepackage{geometry}
\usepackage{fancyhdr}
\usepackage{titling}
\usepackage{booktabs}
\usepackage{hyperref}
\usepackage{xcolor}
\usepackage{longtable}

\geometry{left=2.5cm,right=2.5cm,top=2.5cm,bottom=2.5cm}

\pagestyle{fancy}
\fancyhead[L]{\textit{论文创作过程记录}}
\fancyhead[R]{2026-06-10}
\fancyfoot[C]{\thepage}

\title{论文创作过程记录}
\author{ARS Pipeline v3.12.0}
\date{2026-06-10}

\begin{document}

\maketitle
\tableofcontents
\newpage

\section{论文基本信息}

\begin{tabular}{ll}
\hline
\textbf{项目} & \textbf{内容} \\
\hline
标题 & 技术在场、服务缺席：数智养老服务中数字排斥的生成逻辑与治理路径 \\
领域 & 老龄化/公共管理交叉 \\
字数 & 约 13,135 中文 \\
参考文献 & 18 篇（7 中文 + 11 英文/法规） \\
终稿文件 & 终稿.md + 终稿.docx \\
管线阶段 & 6 阶段（2.5$\rightarrow$3$\rightarrow$4$\rightarrow$3'$\rightarrow$4.5$\rightarrow$5） \\
\hline
\end{tabular}

\section{阶段逐项记录}

\subsection{Stage 2.5 — 完整性审核（Pre-Review Integrity）}

\textbf{输入}：原始初稿（$\sim$8,846 中文字，10 条参考文献）

\textbf{审核结果}：\textbf{PASS WITH NOTES}

\begin{tabular}{lr}
\hline
\textbf{阶段} & \textbf{结果} \\
\hline
Phase A 参考文献 100\% 验证 & 8/10 VERIFIED, 1 NOT FOUND, 1 页码不符 \\
Phase B 引文上下文 30\% 抽检 & 5/6 准确 \\
Phase C 统计数据验证 & 2/10 可独立确认 \\
Phase D 原创性 30\% 抽检 & 85-90\% 原创综合，无抄袭 \\
Phase E 声明验证 30\% 抽检 & 6/13 确认 \\
\hline
\end{tabular}

\textbf{关键发现}：参考文献 4（曹科岩, 方之瑜, 2026）在 WebSearch 中未找到，经本地知识库确认存在原文 PDF 和精读报告，属于 2026 年新刊网络索引延迟。

\textbf{用户决策}：用户要求优先检查本地知识库，确认文献真实存在后进入下一阶段。

\subsection{Stage 3 — 同行评审（Full Review）}

\textbf{配置}：5 位审稿人（主编、方法论专家、老龄化专家、跨学科专家、魔鬼代言人）

\textbf{评审结果}：\textbf{Major Revision} — 综合评分 25.7/35

\begin{tabular}{lrl}
\hline
\textbf{审稿人} & \textbf{评分} & \textbf{关键意见} \\
\hline
主编 & 27/35 & 问题意识强，认识论定位不清 \\
方法论专家 & 25/35 & 框架方法学不透明 \\
老龄化专家 & 27/35 & 群体分化刻画不足 \\
跨学科专家 & 25/35 & STS 理论对话不足 \\
魔鬼代言人 & 19/35 & 核心论点可证伪性存疑 \\
\hline
\end{tabular}

\textbf{修订路线图}：16 条修订要求（P0: 5, P1: 5, P2: 4, P3: 2）

\subsection{Stage 4 — 修订（Major Revision）}

\textbf{修订范围}：16 条全部逐条回应

\begin{tabular}{ll}
\hline
\textbf{优先级} & \textbf{关键变更} \\
\hline
P0 & 声明为概念性论文；全篇"揭示"$\rightarrow$"提出"；转向叙事改为共存叙事；新增理论综合方法说明 \\
P1 & 新增老年群体内部分化；AIQ 操作化定义；补充 7 条国际文献；交叉映射分析 \\
P2 & 需求侧视角；国际比较（日本+欧盟）；政治经济学分析；算法问责可操作化 \\
P3 & 段落拆分；术语统一 \\
\hline
\end{tabular}

\textbf{修订后字数}：$\sim$12,826 中文字（+45\%）

\subsection{Stage 3' — 再审（Verification Review）}

\textbf{评审结果}：\textbf{MINOR REVISION}

\begin{tabular}{lrr}
\hline
\textbf{状态} & \textbf{数量} & \textbf{占比} \\
\hline
PASS & 15 & 93.75\% \\
PARTIAL & 1 & 6.25\% \\
\hline
\end{tabular}

\textbf{剩余问题}：隐性排斥操作化定义承诺未写入正文、R\&R 答复信位置描述不准确。两项均已修复。

\subsection{Stage 4.5 — 最终完整性审核（Final Check）}

\textbf{初始裁决}：FAIL（检出 6 项问题）

\begin{tabular}{lll}
\hline
\textbf{级别} & \textbf{问题} & \textbf{修复} \\
\hline
CRITICAL & 幽灵引用 [9][10] & 已删除 \\
HIGH & 三个正文引用无对应条目 & 已补入 \\
HIGH & 页码错误未修正 & 已修正 \\
MEDIUM & 42\%无来源 & 已删除 \\
MEDIUM & 数据归属不严谨 & 已补充 \\
LOW & 编者笔误 & 已修正 \\
\hline
\end{tabular}

\textbf{最终裁决}：\textbf{PASS}（零问题）

\subsection{Stage 5 — 定稿（Finalize）}

\textbf{格式转换}：
\begin{itemize}
    \item .md 终稿：13,135 中文字，18 条参考文献
    \item .docx 终稿：Pandoc 转换
\end{itemize}

\section{管线统计数据}

\begin{tabular}{lr}
\hline
\textbf{指标} & \textbf{数值} \\
\hline
管线阶段数 & 6 \\
完整性审核轮次 & 2 \\
审稿轮次 & 2 \\
修订条目 & 16 \\
参考文献增长 & 10 $\rightarrow$ 18（+80\%） \\
正文字数增长 & 8,846 $\rightarrow$ 13,135（+48\%） \\
\hline
\end{tabular}

\section{合作质量评估}

\begin{tabular}{lr}
\hline
\textbf{维度} & \textbf{得分} \\
\hline
方向设定能力 & 82/100 \\
智力贡献 & 75/100 \\
质量把关能力 & 85/100 \\
迭代自律 & 78/100 \\
授权效率 & 70/100 \\
元学习 & 70/100 \\
\hline
\textbf{总分} & \textbf{78/100} \\
\hline
\end{tabular}

\textbf{整体评价}：用户展现了良好的学术判断力和质量意识，在关键节点做出了正确的方向性决策。

\textbf{做得好的方面}：坚持从原始初稿出发确保审稿链条完整；对字数的明确要求避免了输出缩水；走完完整流程保证了质量。

\textbf{AI 价值加成}：5 位审稿人配置和魔鬼代言人机制提供了多角度深度反馈；两轮完整性审核通过自零验证有效捕捉幽灵引用和书目错误。

\section{AI 自我反思报告}

\textbf{行为总结}：本轮管线运行整体表现良好，各阶段任务按协议执行。

\textbf{谄媚风险评估}：低（无连续让步，无健康检查警报）

\textbf{AI 在本轮中的不足}：
\begin{enumerate}
    \item Stage 3' 再审结果呈现不完整，先给出汇总表而非完整报告
    \item R\&R 答复信承诺的操作化定义段未实际写入正文
    \item 6 项问题在 Stage 4.5 才暴露，说明中间质量检查不够主动
\end{enumerate}

\textbf{研究失败模式审计日志}：

\begin{tabular}{lll}
\hline
\textbf{模式} & \textbf{最终状态} & \textbf{历史} \\
\hline
幻觉/幽灵引用 & CLEAR & Stage 4.5 清除 \\
书目不准确 & CLEAR & Stage 4.5 修复 \\
不可验证统计数据 & OVERRIDDEN & 42\%删除，其余确认 \\
装饰性引用 & CLEAR & 已删除 \\
过度依赖单一来源 & OVERRIDDEN & 新增 8 条国际文献 \\
抄袭/改写不足 & CLEAR & 全部原创综合 \\
数据造假 & CLEAR & 本地知识库确认 \\
\hline
\end{tabular}

\vfill
\noindent\textit{本报告由 AI 自动生成，反映本轮管线运行情况。}

\noindent\textit{记录日期：2026-06-10}

\end{document}
